\documentclass[11pt]{amsart}
\usepackage{geometry}                % See geometry.pdf to learn the layout options. There are lots.
\geometry{letterpaper}                   % ... or a4paper or a5paper or ... 
%\geometry{landscape}                % Activate for for rotated page geometry
%\usepackage[parfill]{parskip}    % Activate to begin paragraphs with an empty line rather than an indent
\usepackage{graphicx}
\usepackage{amssymb}
\usepackage{epstopdf}
\DeclareGraphicsRule{.tif}{png}{.png}{`convert #1 `dirname #1`/`basename #1 .tif`.png}

\title{EJEMPLO 3.4}
%\date{}                                           % Activate to display a given date or no date

\begin{document}
\maketitle
%\section{}
%\subsection{}
\noindent Supongamos que debemos obtener la suma de los gastos que hicimos en nuestro último viaje, pero no sabemos exactamente cuántos fueron. Los datos son expresados en forma:\\

\noindent Datos: GASTO1, GASTO2 ,...., -1\\
\noindent Donde: \\

\noindent \quad GASTOi \quad  es una variable de tipo real que representa el gasto número i.\\

\noindent \textbf {Explicación de las variables} \\

\noindent SUMGAS: Es una variable de tipo real.
Es un acumulador.
Acumula los datos efectuados.\\

\noindent GASTO: Es una variable de tipo real. Su valor en la primera lectura debe ser verdadero, es decir distinto de -1 . Su valor se modifica en cada vuelta del ciclo. Cuando gasto tome el valor de -1 , entonces el ciclo se detendrá.\\

\noindent En la siguiente tabla podemos observar el seguimiento del algoritmo para los siguientes datos (GASTO): \$2,528, \$3,500, \$1,600, \$1,850, \$150, -1\\
\begin{table}[htbp]
\centering
\label{tab:corrida}
\begin{tabular}{|c|c|c|c|c|c|}
\hline
\multicolumn{3}{ |c| }{\textbf{Tabla 3.3}} \\ \hline
 & \textbf{SUMGAS} & \textbf{GASTO} \\ \hline
Inicia el ciclo $ \rightarrow $ & 0 & 2528 \\ \hline
& 2528 & 3500\\ \hline
& 6028 & 1600 \\ \hline
& 7628 & 1850 \\ \hline
& 9478 & 150\\ \hline
Finaliza el ciclo $ \rightarrow $ & 9628 & -1\\ \hline
\end{tabular}
\end{table}

\end{document}  